\documentclass[11pt,a4paper]{article}

\usepackage[utf8]{inputenc}
\usepackage[T1]{fontenc}
\usepackage[spanish]{babel}
\usepackage{amsmath,amssymb}
\usepackage[margin=2.5cm]{geometry}
\usepackage{booktabs}
\usepackage{tabularx}
\usepackage{hyperref}

% \code{} admite guiones bajos y signos de pesos sin escaparlos
\DeclareRobustCommand{\code}[1]{\texttt{\detokenize{#1}}}
\newcolumntype{L}{>{\raggedright\arraybackslash}X}
\emergencystretch=2em

\title{SpatFD v2: estructura del paquete\\[0.3em]\large Guía para desarrolladores}
\author{Documento de desarrollo}
\date{Octubre de 2026}

\begin{document}
\sloppy
\maketitle

\noindent
Este documento describe cómo está construido SpatFD v2 por dentro: qué hay en
cada archivo, cómo es el objeto \code{SpatFD}, por dónde pasa una predicción y
qué hay que tocar para extender el paquete. Está pensado para quien vaya a
mantenerlo. La teoría de cada método está en la documentación de las funciones
(sección \emph{Theory}) y el diseño del \emph{solver} en
\code{dev/diseno_solver.pdf}; aquí solo se repite lo necesario para entender el
código.

\tableofcontents

% ======================================================================
\section{Principios de diseño}

El paquete sigue la convención de los paquetes de modelado de R (\code{lm},
\code{predict}, \code{update}). Cinco reglas explican casi todas las decisiones
del código:

\begin{enumerate}
  \item \textbf{Los objetos son listas normales con clase.} No hay
        \emph{environments} ni métodos \code{$} sobrescritos. R no copia los
        datos al copiar una lista mientras no se modifiquen, así que compartir
        referencias no ahorra memoria de forma apreciable.
  \item \textbf{Nada se modifica en el sitio.} Toda función que cambia un
        objeto lo devuelve: \code{sfd <- add_model(sfd, ...)}.
  \item \textbf{Lo funcional está oculto.} El usuario solo ve matrices y
        tablas. Las curvas suavizadas, los componentes principales, los datos
        originales y los modelos compilados viven en un único atributo,
        \code{internal}, al que el resto del código accede solo por medio de
        accesores (sección~\ref{sec:accesores}).
  \item \textbf{El ajuste de variogramas está fuera del paquete.} SpatFD
        exporta los \emph{scores} (\code{togstat()}), recibe los modelos
        ajustados (\code{toSpatFD()}, \code{add_model()}) y resuelve la
        predicción con código propio.
  \item \textbf{Mensajes al usuario en inglés, comentarios en español.} Los
        mensajes de error y la documentación roxygen están en inglés; los
        comentarios del código, en español, con un título en mayúsculas al inicio
        de cada bloque (por ejemplo, 	exttt{COMPILACIÓN:}).
\end{enumerate}

% ======================================================================
\section{Flujo de trabajo}

\begin{center}
\begin{tabularx}{\textwidth}{lll L}
\toprule
Paso & Función & Archivo & Resultado \\
\midrule
1 & \code{spatfd()} & \code{spatfd.R} & Objeto con una variable: suavizamiento + PCA funcional \\
2 & \code{+} & \code{spatfd.R} & Objeto con varias variables \\
3 & \code{togstat()} & \code{gstat.R} & Objeto \code{gstat} con un campo por \emph{score} \\
4 & (fuera del paquete) & --- & Variogramas ajustados con \code{gstat} \\
5 & \code{add_model()} & \code{add_model.R} & El objeto con el modelo compilado y registrado \\
6 & \code{predict()} & \code{predict.R} & Objeto \code{SpatFDpred}: curvas y varianzas \\
7 & \code{plot()}, \code{print()} & \code{methods-*.R} & Gráficos y resúmenes \\
8 & \code{leave_one_out()} & \code{leave_one_out.R} & Objeto \code{SpatFDloo}: comparación de modelos y métodos \\
\bottomrule
\end{tabularx}
\end{center}

% ======================================================================
\section{Organización de los archivos}

Cada archivo agrupa una función exportada (o una familia) con sus funciones
internas. Las funciones internas que usan varios archivos están en
\code{utils.R}. La tabla lista todo lo que hay en \code{R/}.

\begin{center}
\small
\begin{tabularx}{\textwidth}{l L L}
\toprule
Archivo & Exportado & Interno \\
\midrule
\code{SpatFD-package.R} & Documentación del paquete & Todos los \code{@importFrom} \\
\code{spatfd.R} & \code{spatfd()}, \code{+.SpatFD} & --- \\
\code{methods-SpatFD.R} & \code{print}, \code{summary}, \code{plot} de \code{SpatFD}; \code{var_names()} & \code{print.summary.SpatFD}, \code{plot_components()} \\
\code{gstat.R} & \code{togstat()}, \code{toSpatFD()} y sus métodos & --- \\
\code{add_model.R} & \code{add_model()} & \code{compile_model()} \\
\code{models.R} & \code{models}, \code{availablemodels()} & \code{check_params()} \\
\code{predict.R} & \code{predict.SpatFD} & \code{predict_scores()}, \code{predict_lambda()}, \code{prediction_setup()}, \code{krige_scores()} \\
\code{kriging.R} & --- & \code{score_covariances()}, \code{cross_dist()}, \code{solve_kriging()}, \code{constrained_solve()} \\
\code{methods-SpatFDpred.R} & \code{print}, \code{plot} de \code{SpatFDpred} & \code{predicted_fd()}, \code{observed_fd()}, \code{map_grid()}, \code{plot_curves()}, \code{plot_maps()} \\
\code{leave_one_out.R} & \code{leave_one_out()}, \code{print.SpatFDloo} & \code{curve_distance()}, \code{refit()} \\
\code{utils.R} & --- & Constructor, accesores, \code{score_ids()}, \code{model_table()}, \code{models_message()}, \code{evaluate_curves()}, \code{inner_product()} \\
\code{data.R} & Documentación de \code{air_mexico}, \code{air_bogota}, \code{vowels} & --- \\
\bottomrule
\end{tabularx}
\end{center}

Otras carpetas:
\begin{itemize}
  \item \code{tests/testthat/}: una prueba por archivo de \code{R/}
        (sección~\ref{sec:pruebas}).
  \item \code{dev/}: scripts de uso (\code{01_uso_paquete.R}) y de comparación
        con \code{gstat} (\code{02_comparacion_gstat.R}), y documentos técnicos.
        Está en \code{.Rbuildignore}.
  \item \code{archive/}: versiones anteriores y material de referencia. Está en
        \code{.gitignore}.
  \item \code{man/} y \code{NAMESPACE}: generados por roxygen. No se editan a
        mano; se regeneran con \code{devtools::document()}.
\end{itemize}

% ======================================================================
\section{El objeto \code{SpatFD}}
\label{sec:objeto}

\subsection{Estructura}

\begin{footnotesize}
\begin{verbatim}
sfd                                   lista, clase "SpatFD"
|-- $scores     matriz: un sitio por fila; coordenadas y luego <var>.PC1, <var>.PC2, ...
|-- $variance   data.frame: una fila por variable (basis, no. PC, variance)
|-- $models     data.frame: una fila por modelo registrado (tabla resumen)
`-- attr(, "internal")                lista
    |-- $coords     matriz con las coordenadas de todos los sitios
    |-- $variables  lista, una entrada por variable:
    |   `-- $PM10
    |       |-- $data   matriz original (tiempos x sitios)
    |       |-- $fd     curvas suavizadas (objeto fd de fda)
    |       |-- $pca    resultado de fda::pca.fd()
    |       `-- $args   list(t, basis, nbasis, lambda): para repetir el análisis
    `-- $models     lista, una entrada por modelo registrado:
        `-- $lmc
            |-- $table       data.frame del modelo (var, model, rango, sigma, alpha1)
            |-- $variograms  lista de funciones gamma(h), una por id de variograma
            `-- $sills       vector con la meseta total de cada id
\end{verbatim}
\end{footnotesize}

\subsection{Invariantes}

Todo el código supone que se cumple lo siguiente. Cualquier función que cree o
modifique un objeto debe mantenerlo.

\begin{itemize}
  \item Las filas de \code{$scores} son todos los sitios del objeto, ordenados
        por nombre. Un sitio donde no se observó una variable tiene \code{NA}
        en sus \emph{scores}.
  \item Las columnas de \code{$scores} son primero las coordenadas (las de
        \code{internal$coords}) y luego los \emph{scores}, con nombres
        \code{<variable>.PC<k>}. El orden de las columnas importa: es el que
        recibe \code{gstat} (sección~\ref{sec:gstat}).
  \item Los \emph{ids} de los variogramas son los nombres de los
        \emph{scores} (\code{PM10.PC1}) para los directos, y dos de ellos
        unidos por un punto (\code{NO2.PC1.PM10.PC1}) para los cruzados, en
        cualquier orden.
  \item \code{$models} es siempre \code{model_table()} del objeto. Solo
        \code{new_spatfd()} y \code{set_models()} la escriben.
  \item \code{var_names()} devuelve las variables en orden alfabético. Las
        funciones que recorren variables usan ese orden.
\end{itemize}

\subsection{Constructor, accesores y modificadores}
\label{sec:accesores}

Están en \code{utils.R}. El resto del paquete \textbf{nunca} usa
\code{attr(x, "internal")} directamente; así, si la estructura interna
cambia, solo cambia este archivo.

\begin{center}
\small
\begin{tabularx}{\textwidth}{l L}
\toprule
Función & Qué hace \\
\midrule
\code{new_spatfd(scores, variance, coords, variables, models)} & Arma el objeto y calcula \code{$models} \\
\code{get_coords(sfd)} & Coordenadas de todos los sitios \\
\code{get_variable(sfd, v)} & Entrada completa de la variable: \code{data}, \code{fd}, \code{pca}, \code{args} \\
\code{get_fd(sfd, v)}, \code{get_pca(sfd, v)} & Curvas suavizadas y PCA de la variable \\
\code{get_models(sfd)} & Lista de modelos compilados \\
\code{set_pca(sfd, v, pca)} & Devuelve el objeto con el PCA de \code{v} reemplazado \\
\code{set_models(sfd, models)} & Devuelve el objeto con los modelos reemplazados y \code{$models} recalculada \\
\code{score_ids(sfd)} & Vector con la variable de cada \emph{score}, nombrado con el id del \emph{score} \\
\code{model_table(sfd)} & La tabla resumen de los modelos \\
\code{models_message(sfd)} & La tabla resumen como texto, para los mensajes de error \\
\bottomrule
\end{tabularx}
\end{center}

\subsection{Construcción: \code{spatfd()}}

\begin{enumerate}
  \item Verifica los argumentos con \code{stopifnot()} y mensajes con nombre.
  \item Suaviza: base B-splines o Fourier sobre \code{range(t)}, penalización
        \code{vec2Lfd()} y \code{smooth.basis()}.
  \item Calcula el PCA con \code{pca.fd(fd, nharm)}. Si se da \code{var},
        repite el PCA con el mínimo número de componentes que alcanza esa
        varianza.
  \item Arma \code{scores}, \code{variance} y la entrada de la variable en
        \code{variables}, con \code{args} para poder repetir el análisis, y
        llama a \code{new_spatfd()}.
\end{enumerate}

\subsection{Combinación: \code{+.SpatFD}}

\begin{itemize}
  \item Exige que las coordenadas tengan las mismas columnas.
  \item Si una variable está en los dos objetos, conserva la del primero y
        avisa con un \emph{warning}.
  \item Coordenadas: unión de los sitios, primero los del primer objeto.
  \item \emph{Scores}: unión completa por sitio (\emph{full join}) con
        \code{match()}. No se indexan filas de un \emph{data frame} por nombre,
        porque R hace coincidencias parciales: el sitio \code{s1} tomaría los
        \emph{scores} de \code{s10}. Hay una prueba para este caso.
  \item \code{variables}: concatenación de las listas.
  \item Modelos: los del primer objeto y luego los del segundo. Un nombre
        repetido se renombra en secuencia (\code{model 2}, \dots) con un
        mensaje.
\end{itemize}

% ======================================================================
\section{Intercambio con \code{gstat}}
\label{sec:gstat}

\subsection{\code{togstat()}}

Crea un objeto \code{gstat} con un campo por \emph{score} (fórmula
\code{score ~ 1}) y quita en cada campo los sitios sin observación. Los campos
siguen el orden de las columnas de \code{$scores}. \textbf{Ese orden no es
indiferente}: \code{gstat::fit.lmc()} llega a ajustes distintos según el orden
de los campos. Hubo una regresión por ordenar alfabéticamente, y hay una prueba
en \code{test-spatfd.R} que lo impide.

\subsection{\code{toSpatFD()}: el esquema tabular}

Todo modelo, venga de donde venga, se convierte a un \emph{data frame} con una
fila por estructura y las columnas:

\begin{center}
\begin{tabular}{ll}
\toprule
Columna & Contenido \\
\midrule
\code{var} & id del variograma (directo o cruzado) \\
\code{model} & nombre del modelo en la lista \code{models} \\
\code{rango} & parámetro de rango \\
\code{sigma} & meseta parcial (puede ser negativa en los cruzados) \\
\code{alpha1} & parámetro de forma (\code{kappa} en \code{gstat}) \\
\bottomrule
\end{tabular}
\end{center}

\code{toSpatFD()} es genérica. Métodos actuales: \code{gstat} (resultado de
\code{fit.lmc()}), \code{list} (lista nombrada de \code{variogramModel}, que se
envuelve como un \code{gstat}) y \code{default} (error). La traducción de
nombres de \code{gstat} está en \code{toSpatFD.gstat()}. La anisotropía se
descarta.

% ======================================================================
\section{Modelos de variograma}

\subsection{La lista \code{models}}

\code{models} (exportada) es una lista nombrada de funciones
\code{function(h, rango, sigma, alpha1, ...)} que devuelven $\gamma(h)$. Hay 18
espejos de \code{gstat}, que reproducen sus valores (prueba en
\code{test-models.R}), y 4 modelos propios. Todas validan sus parámetros con
\code{check_params()}, que define los dominios:

\begin{center}
\begin{tabular}{ll}
\toprule
Dominio & Usado en \\
\midrule
\code{'positive'} (por defecto) & \code{rango} y \code{alpha1} de casi todos \\
\code{'non-negative'} & \code{rango} de \code{linear} (0 = sin meseta) \\
\code{'up to 2'} & \code{rango} de \code{power}, \code{alpha1} de \code{stable} \\
\bottomrule
\end{tabular}
\end{center}

\code{check_params()} solo revisa los parámetros que recibe: \code{nugget} no
le pasa \code{rango}. El signo de \code{sigma} no se revisa.

\subsection{Registro: \code{add_model()}}

\begin{enumerate}
  \item Si el modelo no viene en el esquema tabular, lo traduce con
        \code{toSpatFD()}.
  \item Verifica que los ids sean \emph{scores} del objeto o parejas de
        ellos, y que los modelos existan en \code{models}.
  \item Compila con \code{compile_model()} y calcula las mesetas totales por id.
  \item Guarda la entrada (\code{table}, \code{variograms}, \code{sills}) con
        \code{set_models()}. Un nombre repetido reemplaza el modelo sin cambiar
        su posición; sin nombre, se usa \code{next_model_name()}.
\end{enumerate}

\subsection{Compilación: \code{compile_model()}}

Para cada id, cada fila de la tabla se convierte en un \emph{closure}
\code{function(h) do.call(model, c(list(h), args))}, y el variograma es
la suma de esos \emph{closures}. Los parámetros en \code{NA} se omiten, de modo
que el modelo usa su valor por defecto. Cada estructura se evalúa una vez con
\code{numeric(0)} al compilar: así un parámetro inválido falla al registrar el
modelo y no al predecir.

% ======================================================================
\section{Predicción}

\subsection{Recorrido de una llamada}

\begin{footnotesize}
\begin{verbatim}
predict.SpatFD(object, newcoords, model, variable, method, kriging, n)
 |-- si newcoords es sf/sfc:  map_grid()  -> grilla regular dentro del mapa
 |-- method = "scores":  predict_scores()
 |     `-- krige_scores()
 |           |-- prediction_setup()     verifica y prepara
 |           |-- score_covariances()    C(h) = meseta - gamma(h), por pareja de scores
 |           `-- solve_kriging()        resuelve el sistema (simple u ordinario)
 |-- method = "lambda":  predict_lambda()
 |     |-- prediction_setup()
 |     |-- score_covariances()
 |     `-- sistema propio (bloques por variable), con constrained_solve() si es ordinario
 `-- arma el objeto SpatFDpred
\end{verbatim}
\end{footnotesize}

\subsection{\code{prediction_setup()}}

Es el paso común a los dos métodos. Escoge el modelo (si hay uno solo, no hace
falta nombrarlo), elige las variables (por defecto, todas las que cubre el
modelo), ordena las columnas de \code{newcoords} según las del objeto (por
nombre, o por posición si no tienen nombre) y separa, para cada \emph{score},
los sitios donde se observó y sus valores. Devuelve:

\begin{center}
\begin{tabular}{ll}
\toprule
Elemento & Contenido \\
\midrule
\code{model}, \code{compiled} & Nombre del modelo y su entrada compilada \\
\code{pcs} & Lista: los ids de los \emph{scores} de cada variable \\
\code{sites}, \code{values} & Sitios y valores observados de cada \emph{score} \\
\code{newcoords} & Coordenadas verificadas, con nombres de fila \\
\bottomrule
\end{tabular}
\end{center}

Todos los errores sobre el modelo incluyen \code{models_message()}, la tabla de
modelos del objeto.

\subsection{\code{score_covariances()}}

Devuelve una matriz-lista $q \times q$ (una fila y una columna por
\emph{score}) con la función de covarianza de cada pareja, $C(h) = C(0) -
\gamma(h)$. Las parejas sin variograma cruzado quedan en \code{NULL}, es decir,
independientes. Los modelos sin meseta (\code{unbounded_models}: \code{linear},
\code{power}, \code{logarithmic}, \code{spline}, \code{legendre}) solo se
admiten con kriging ordinario, con la pseudo-covarianza $C(h) = -\gamma(h)$
(sección~\ref{sec:ordinario}); con kriging simple dan un error.

\subsection{Kriging simple y ordinario}
\label{sec:ordinario}

\code{predict()} y \code{leave_one_out()} reciben \code{kriging}, por defecto
\code{"ordinary"}. El simple supone media conocida e igual a cero (los
\emph{scores} están centrados). El ordinario supone una media constante
desconocida y agrega una restricción por campo: los pesos del campo que se
predice suman 1 y los de cada uno de los demás suman 0 (cokriging ordinario).
\code{constrained_solve()} arma y resuelve, por LU, el sistema
\[
  \begin{pmatrix} \Sigma & F \\ F^T & 0 \end{pmatrix}
  \begin{pmatrix} W \\ M \end{pmatrix} =
  \begin{pmatrix} C_0 \\ E \end{pmatrix},
  \qquad
  \sigma^2 = C(0) - \textstyle\sum_{\text{columna}} [C_0; E] \circ [W; M],
\]
donde $F$ marca las filas de cada campo y $E$ el campo objetivo de cada
columna. Dos consecuencias de las restricciones guían el diseño:
\begin{itemize}
  \item \textbf{Las constantes se cancelan.} Sumar una constante a un bloque de
        covarianzas no cambia ni los pesos ni los multiplicadores ni la
        varianza. Por eso los modelos sin meseta funcionan con $C(h) =
        -\gamma(h)$, y los que tienen meseta conservan $C(h) = C(0) - \gamma(h)$,
        que está mejor condicionada.
  \item \textbf{Da igual centrar las curvas.} En lambda, como los pesos de la
        variable objetivo suman 1 y los de las demás 0, las medias se cancelan:
        $\hat\chi^r = \sum_p \sum_i \lambda^{rp}_i \chi^p_{s_i}$ con las curvas
        sin centrar. El código sigue sumando la media a las curvas centradas,
        que es la misma cuenta y vale también para el simple.
\end{itemize}
La matriz del ordinario no es definida positiva, así que no se usa Cholesky.
Las dos variantes se comparan con \code{gstat} (con y sin \code{beta = 0}) en
\code{test-predict.R}, \code{test-ordinary.R} y \code{dev/02_comparacion_gstat.R}.

\subsection{Método \code{"scores"}: \code{solve_kriging()}}

Resuelve el kriging de $q$ campos escalares (los \emph{scores}). El
detalle matemático está en \code{dev/diseno_solver.pdf}; las decisiones de
implementación son:

\begin{itemize}
  \item \textbf{Grupos.} Los campos conectados, directa o indirectamente, por
        covarianzas cruzadas forman un grupo, y cada grupo se resuelve por
        separado. Sin cruzados hay un kriging por \emph{score}; con un modelo
        lineal de corregionalización completo hay un solo sistema de cokriging.
  \item \textbf{Distancias.} Se calculan una vez por conjunto distinto de
        sitios (los \emph{scores} de una variable comparten sitios) y se
        guardan en caché.
  \item \textbf{Simple: un Cholesky por grupo.} Con $\Sigma = R^T R$ y $V =
        R^{-T} C_0$, la predicción es $V^T (R^{-T} z)$ y la varianza $C(0) -
        \sum V^2$ por columna, sin calcular los pesos. Si Cholesky falla, la
        matriz no es definida positiva: hay sitios duplicados o el modelo no es
        válido.
  \item \textbf{Ordinario: un LU por grupo} con \code{constrained_solve()}, con
        una restricción por campo del grupo.
  \item \textbf{Distancia cero exacta.} \code{cross_dist()} calcula la
        distancia coordenada a coordenada para que la de un sitio consigo
        mismo sea exactamente 0. De eso depende el efecto pepita.
\end{itemize}

\code{predict_scores()} reconstruye después cada curva sobre los coeficientes
de la base: $\text{coefs} = \Xi\, \hat{f} + \text{coefs}(\mu)$, con $\Xi$ los
coeficientes de las funciones propias.

\subsection{Método \code{"lambda"}: cokriging funcional}

\code{predict_lambda()} implementa el predictor de Bohórquez, Giraldo y Mateu
(2017), con las curvas centradas:
\[
  \hat\chi^r_{s_0}(t) = \mu^r(t) + \sum_{p} \sum_{i}
  \lambda^{rp}_i \bigl(\chi^p_{s_i}(t) - \mu^p(t)\bigr).
\]
Pasos:

\begin{enumerate}
  \item \textbf{Intervalo común} $T$: la intersección de los intervalos de las
        variables. Si no coinciden, avisa con un mensaje; si no se cruzan, da
        un error.
  \item \textbf{Matrices} $G^{pq}_{kl} = \int_T \xi^p_k \xi^q_l\,dt$. Para una
        variable observada en todo $T$, $G^{pp}$ es exactamente la identidad
        (no se calcula). En los demás casos se usa \code{inner_product()}.
  \item \textbf{Covarianza entre curvas} de cada pareja de variables,
        $\Omega_{pq}(h) = \sum_{k,l} G^{pq}_{kl}\, C_{pk,ql}(h)$, como función de
        $h$. Es \code{NULL} si no hay ningún término.
  \item \textbf{Sistema} con un bloque por pareja de variables. Simple:
        $\Sigma \lambda = C_0$ con un Cholesky. Ordinario:
        \code{constrained_solve()} con una restricción por variable. Para cada
        variable objetivo solo cambian $C_0$ y la restricción.
  \item \textbf{Curvas}: se evalúan en los tiempos de la variable dentro de
        $T$. Esa matriz es exactamente la predicción y es la que recibe el
        usuario.
  \item \textbf{Varianza}: $\Omega_{rr}(0) - C_0^T \lambda$ (menos el
        multiplicador en el ordinario), con las mesetas del modelo. No puede
        ser negativa.
  \item \textbf{Versión funcional} (solo para graficar): la matriz se proyecta
        por mínimos cuadrados en la base de la variable, o en B-splines sobre
        $T$ si el intervalo es más corto.
\end{enumerate}

Este método no usa \code{solve_kriging()}: el sistema es pequeño (un bloque por
variable, no por \emph{score}) y escribirlo aparte resultó más claro que
generalizar el \emph{solver} a valores matriciales.

\subsection{\code{inner_product()}}

\code{fda::inprod()} se detiene en un error relativo de $10^{-4}$, que no basta
para un sistema que se verifica a $10^{-8}$. \code{inner_product()} usa la
matriz exacta de productos internos de la base (\code{eval.penalty}) cuando las
dos curvas están en la misma base y se integra en todo su intervalo. Si no,
integra por Gauss-Legendre de 5 puntos en 1000 tramos (error de $10^{-15}$ con
Fourier y $10^{-9}$ con B-splines en las pruebas).

\subsection{El objeto \code{SpatFDpred}}

\begin{footnotesize}
\begin{verbatim}
pred                                  lista, clase "SpatFDpred"
|-- $PM10
|   |-- $pred   matriz: tiempos x sitios nuevos (curvas evaluadas)
|   `-- $var    vector: varianza de predicción por sitio nuevo
|-- attr(, "newcoords")   coordenadas de los sitios nuevos
|-- attr(, "method")      "scores" o "lambda"
|-- attr(, "crossed")     si se usó alguna covarianza cruzada
|-- attr(, "map")         el mapa sf, si se predijo sobre uno
`-- attr(, "internal")
    |-- $pred       curvas predichas como fd, por variable
    |-- $observed   curvas observadas como fd, por variable
    `-- $coords     coordenadas de los sitios observados
\end{verbatim}
\end{footnotesize}

\noindent
Contrato de los métodos: \code{predict_scores()} y \code{predict_lambda()}
devuelven \code{list(curves, newcoords, crossed)}, donde \code{curves} tiene por
variable \code{fd}, \code{pred} y \code{var}. Un método nuevo solo tiene que
cumplir ese contrato y agregarse al \code{switch} de \code{predict.SpatFD}.

% ======================================================================
\section{Gráficos}

\begin{center}
\small
\begin{tabularx}{\textwidth}{l L}
\toprule
Función & Qué dibuja \\
\midrule
\code{plot.SpatFD}, \code{mode = "functional"} & Curvas suavizadas, con \code{plot.fd} de \code{fda} \\
\code{plot.SpatFD}, \code{mode = "pca"} & Todos los componentes en una página (\code{plot_components()}), con \code{plot.pca.fd} componente a componente \\
\code{plot.SpatFDpred}, \code{mode = "curves"} & Observadas en gris y predichas en color (\code{plot_curves()}) \\
\code{plot.SpatFDpred}, \code{mode = "map"} & Predicción en un instante y varianza, lado a lado (\code{plot_maps()}) \\
\code{availablemodels()} & Los modelos de variograma en una grilla de $5 \times 5$ \\
\bottomrule
\end{tabularx}
\end{center}

\textbf{Política de \code{ask}.} Ningún gráfico del paquete pide
\code{<Enter>}. Cada uno guarda \code{par(ask = FALSE)} al empezar y restaura el
estado del dispositivo al salir con \code{on.exit(..., add = TRUE)}. La razón:
\code{fda::plot.pca.fd()} dibuja una página por componente y activa
\code{ask}; si el dispositivo queda en ese modo, R pide \code{<Enter>} antes de
cada gráfico nuevo, aunque sea el único. La única excepción es
\code{plot(sfd, multiple = TRUE)}, que el usuario pide explícitamente. La prueba
\code{test-plots.R} deja el dispositivo en modo \code{ask} y comprueba con el
\emph{hook} \code{before.plot.new} que ninguna página nueva se abre en ese modo.

\textbf{Mapas.} \code{map_grid()} toma cerca de \code{n} puntos regulares
dentro del mapa con \code{sf::st_sample(type = "regular")}. \code{plot_maps()}
reconstruye la grilla para \code{image()} a partir del paso mínimo entre
coordenadas y arma la página con \code{layout()}: título, subtítulo, dos barras
de color y dos mapas.

% ======================================================================
\section{Validación cruzada: \code{leave_one_out()}}

Para cada sitio de la variable:

\begin{enumerate}
  \item \code{refit()} repite el suavizamiento y el PCA de la variable sin el
        sitio, con los mismos \code{args} y el mismo número de componentes. Las
        otras variables se reconstruyen una sola vez, sin quitar nada.
  \item Las funciones propias nuevas toman el signo de las originales (por el
        signo de su producto interno). El signo de una componente es arbitrario,
        pero los variogramas cruzados del modelo dependen de él.
  \item Se registran de nuevo los modelos con su tabla original: los
        variogramas \textbf{no} se vuelven a ajustar.
  \item Se predice el sitio con cada combinación de modelo y método, y se mide
        el error con \code{curve_distance()}: la norma $L^2$ de la diferencia,
        exacta si las curvas comparten base y por \code{inner_product()} si no
        (por ejemplo, lambda sobre un intervalo común más corto).
\end{enumerate}

\noindent
Devuelve un \code{SpatFDloo}: \code{summary} (una fila por combinación,
ordenada por error medio), \code{errors} (sitios $\times$ combinaciones) y
\code{pred} (por combinación, matriz tiempos $\times$ sitios; \code{NA} fuera
del intervalo en que se predijo).

% ======================================================================
\section{Pruebas}
\label{sec:pruebas}

Se corren con \code{devtools::test()}. Cada archivo usa datos simulados
pequeños (10 sitios, 50 tiempos) y \code{gstat} como referencia externa
siempre que existe.

\begin{center}
\small
\begin{tabularx}{\textwidth}{l L}
\toprule
Archivo & Qué verifica \\
\midrule
\code{test-spatfd.R} & Construcción, \code{+} (incluido el caso \code{s1}/\code{s10}), orden de \code{togstat()}, que solo se expongan \code{scores}, \code{variance} y \code{models} \\
\code{test-models.R} & Los 18 espejos contra \code{gstat::variogramLine()}, estructuras anidadas, dominios de los parámetros, \code{toSpatFD()} \\
\code{test-predict.R} & Kriging y cokriging de \emph{scores} contra \code{gstat}; lambda univariado contra las fórmulas escritas a mano; cokriging funcional contra el sistema por bloques; intervalos distintos; mapas; validación de argumentos; registro de modelos \\
\code{test-leave_one_out.R} & Cada error es el de repetir el análisis sin el sitio; diseños; modelos multivariados; intervalos distintos \\
\code{test-ordinary.R} & Kriging y cokriging ordinario contra \code{gstat}; modelos sin meseta; lambda ordinario independiente de centrar, contra el sistema con restricciones; \code{leave_one_out()} ordinario por defecto \\ \code{test-plots.R} & Ningún gráfico pide \code{<Enter>}; los componentes van en una página; \code{plot()} devuelve el objeto de forma invisible \\
\bottomrule
\end{tabularx}
\end{center}

\noindent
Las pruebas pueden usar las funciones internas (\code{testthat} carga el
paquete con \code{load_all()}). Además, \code{dev/02_comparacion_gstat.R} repite
las comparaciones con \code{gstat} sobre los datos reales de México y se detiene
si alguna diferencia relativa supera $10^{-8}$.

% ======================================================================
\section{Cómo extender el paquete}

\subsection{Agregar un modelo de variograma}
\begin{enumerate}
  \item Agregar la función a la lista \code{models} en \code{models.R}, con
        \code{check_params()} al inicio y \code{h = 0} resuelto explícitamente
        si la fórmula da \code{NaN} ahí.
  \item Si tiene equivalente en \code{gstat}, agregar la traducción en
        \code{toSpatFD.gstat()} y el caso en la prueba de espejos.
  \item Si no tiene meseta, agregarlo a \code{unbounded_models} en
        \code{kriging.R}: solo se podrá usar con kriging ordinario.
  \item Actualizar la documentación de \code{models}.
\end{enumerate}

\subsection{Agregar un traductor (por ejemplo, \code{geoR})}
Escribir un método \code{toSpatFD.<clase>()} en \code{gstat.R} (o en un archivo
propio) que devuelva el esquema tabular. \code{add_model()} lo usa sin cambios.

\subsection{Agregar un método de predicción}
Escribir \code{predict_<nombre>()} con el contrato de la sección anterior,
agregarlo al \code{switch} y a \code{match.arg()} de \code{predict.SpatFD}, y
documentarlo en \code{@details} y \code{@section Theory}.
\code{leave_one_out()} lo admite si se agrega al diseño por defecto.

\subsection{Cambiar la estructura interna}
Solo hay que tocar \code{utils.R} (constructor y accesores) y las pruebas de
estructura de \code{test-spatfd.R}. Si un cambio obliga a tocar otros archivos,
probablemente se está saltando un accesor.

\subsection{Antes de cada \emph{commit}}
\begin{footnotesize}
\begin{verbatim}
devtools::document()     # NAMESPACE y man/
devtools::test()
source("dev/02_comparacion_gstat.R")
devtools::check()        # 0 errores y 0 warnings
\end{verbatim}
\end{footnotesize}

% ======================================================================
\section{Limitaciones conocidas}

\begin{itemize}
  \item Media constante: el kriging ordinario estima una media constante;
        todavía no hay kriging universal (media con tendencia espacial).
  \item Vecindario global y distancias euclidianas: sirve para coordenadas
        proyectadas y un número moderado de sitios.
  \item El método lambda con varias variables supone que los \code{t} de
        \code{spatfd()} están en una escala común. En los datos de México
        (\code{t = 1, 2, ...}) eso equivale a suponer que PM10 y NO2 empiezan a
        la misma hora.
  \item Solo se traducen modelos de \code{gstat}, sin anisotropía.
  \item Falta migrar desde la versión CRAN el muestreo óptimo, la
        clasificación y la simulación.
\end{itemize}

\end{document}
